O Aprendizado Federado permite treinar modelos em dispositivos de borda sem centralizar dados brutos, mas a heterogeneidade estatística entre clientes (dados não-IID) degrada a convergência e a acurácia do algoritmo padrão, o FedAvg. O FedProx propõe mitigar o problema com um termo proximal na função objetivo local, porém a literatura recente apresenta resultados contraditórios sobre quando ele de fato supera o FedAvg e raramente mede o custo computacional que impõe aos clientes. Este trabalho propõe uma pesquisa experimental que comparará FedAvg e FedProx, tendo SCAFFOLD, FedAvgM e FedAdam como referências, em simulações com o framework Flower sobre CIFAR-10 e Fashion-MNIST. A heterogeneidade será controlada por partição de Dirichlet ($\alpha \in \{0{,}05; 0{,}1; 0{,}5; 1; 100\}$), partição patológica e desbalanceamento de quantidade, e o hiperparâmetro proximal variará em $\mu \in \{0{,}01; 0{,}1; 1\}$, com sementes pareadas (cinco exploratórias e dez confirmatórias), teste de Wilcoxon e teste de equivalência. Serão avaliados acurácia global, estabilidade, rodadas até a acurácia-alvo, dispersão entre clientes e tempo de computação local, este último também em uma CPU restrita. A contribuição esperada é um mapa de trade-off $\alpha \times \mu \times$ (acurácia, custo) que indica em quais níveis e tipos de heterogeneidade o FedProx compensa.

\textcolor{red}{Repositório do projeto: \linkProjeto}
